\section{The boundary and its retained kernel under each original relation}
\label{GraphV3:reader-v3-BRU:sec:bru-original-relations}

This calculation continues the incoming boundary observation A1 and the
boundary/kernel metric A4--A18. It calculates those metrics and every
cross-degree overlap from the original relation rows. All maps below are
on the actual packet, before taking a completion. In particular the
boundary observation does not remove its kernel.

\subsection{The packet, observation, source and finite relation rows}

Retain an actual hypothetical complete-order quartet
\[
 \rho=\tfrac12+\delta+i\gamma,\quad 0<\delta<\tfrac12,\quad\gamma>2,
 \quad h(s)=\prod_{r\in\{\rho,\bar\rho,1-\rho,1-\bar\rho\}}(s-r)^m,
 \quad n=4m,
 \tag{BRU1}
\]
where $m\geq1$ is its complete zero order in $g=2\xi$.
Thus $v_h=g/h$ is the original entire amplitude and
$U=M_{j_hv_h}$ is its invertible multiplication matrix on
$E_h=\mathbb{C}[s]/(h)$. No value of this unit is assigned.
For $k\geq1$ the original cyclic coefficient space and action are
\[
 \begin{gathered}
 \ell_k=1+k(m-1),\qquad q=(k+1)^2\ell_k,\qquad E=\mathbb{C}[S]/(\chi),\\
 \chi(S)=\prod_{a,b=0}^k
 (S-k/2-(2a-k)\delta-i(2b-k)\gamma)^{\ell_k},\qquad
 A=M_S,\\
 \iota:E\longrightarrow E_h^{\otimes k},\qquad
 [P]\longmapsto[P(s_1+\cdots+s_k)],\qquad
 \eta=U^{\otimes k}\iota .
 \end{gathered}\tag{BRU2}
\]
The original cyclic inclusion is injective: on each ordered primary
tensor the sum of the nilpotents has nilpotency index
$1+k(m-1)$, since its power $k(m-1)$ has nonzero coefficient
$(k(m-1))!/((m-1)!)^k$ on the product of the top powers. Its next
power is zero by total degree. All sums in BRU2 are distinct for
different $(a,b)$ because $\delta,\gamma>0$. Each occurs among the
ordered factors, so their product with exactly these exponents is
the minimal polynomial of the tensor sum. This proves the stated
kernel of polynomial evaluation and injectivity of $\iota$.

Fix $u\ne0$ on the original logarithm branch and $t=0$. Let
$\Pi$ be the original period matrix with rays
$\Gamma_j=-\ell_0+\ell_j$, angles
$(\arg u+\pi+2\pi j)/(n+1)$, and entries
$\int_{\Gamma_j}e^{\Phi_h(s)/u}s^b\,ds$, where $\Phi_h'=h$ and
$\Phi_h(0)=0$. Its full determinant and convergence are proved in
PDE1--13 and IPM1--7, included in the source package. Let $C$ have
rows describing $\Gamma_j\mapsto\Gamma_{j+1}-\Gamma_1$ for $j<n$
and $\Gamma_n\mapsto-\Gamma_1$. Set
\[
 P_k=\frac1{n+1}\sum_{a=0}^{n}(C^{-a})^{\otimes k},\qquad
 L=P_k\Pi^{\otimes k}\eta:E\longrightarrow(\mathbb{C}^n)^{\otimes k}.
 \tag{BRU3}
\]
For the algebra used here, $C^{n+1}=I$ follows by applying the
cyclic permutation to the reduced endpoint basis $e_j-e_0$.
Multiplying the finite sum therefore gives $P_k^2=P_k$.
We use the literal coefficient map BRU3, without evaluating a
formal stationary series or identifying its action with Frobenius.
Choose a basis inclusion $\jmath:B=\operatorname{im}L\hookrightarrow
(\mathbb{C}^n)^{\otimes k}$, and define the surjection
$\Lambda:E\to B$ by $\jmath\Lambda=L$. Its kernel is $K$.
Fix a basis inclusion $I:K\hookrightarrow E$, and a coefficient
section $S_*$ with $\Lambda S_*=I_B$. Gaussian elimination on the
specified finite matrix constructs these maps, including at a
rank drop by using its actual image there. Keep these choices fixed
as the source metric varies. Write $r=\dim B$ and $d_K=q-r$.
Zero-dimensional factors use their unique empty inverse and determinant one.

The actual source interpolation is, for $0\leq x\leq1$,
\[
 \begin{gathered}
 S=k/2+iy,\qquad w_h(y)=|v_h(1/2+iy)|^2/(2\pi),\qquad
 m_{h,k}=w_h^{*k},\\
 c_a=2^{1-2a}\Gamma(2a),\qquad
 r_k^\Gamma(y)=\frac{(\sqrt{2\pi})^k}{c_{k/4}}
                   \frac{|\Gamma(k/4+iy/2)|^2}{2\pi},\\
 d\mu_x(y)=((1-x)r_k^\Gamma(y)+xm_{h,k}(y))\,dy,\qquad
 \langle P,Q\rangle_x=\int\overline{P(S)}Q(S)\,d\mu_x(y).
 \end{gathered}\tag{BRU4}
\]
All original masses remain in this measure. Its polynomial moments
are finite by the inherited theta growth and vertical decay estimates.
Its Gram is positive: the Gamma density is positive, and $v_h$ is
a nonzero entire function with isolated zeros, so $w_h$ is positive
almost everywhere; its repeated convolution is positive almost
everywhere. A nonzero polynomial cannot vanish on a set of positive
Lebesgue measure. These statements also prove positivity at $x=0,1$.

Let $J_N:\mathcal P_N\to E$ be the original monic remainder map,
$s:E\to\mathcal P_{q-1}$ its increasing coefficient section, and
$R_N:E\to\mathcal P_N$ its minimum-norm section for $N\geq q-1$.
Explicitly $M_N(x)=(\int(\overline S)^aS^b\,d\mu_x)_{0\leq a,b\leq N}$
is the source Gram in its increasing power frame, and $M_x=M_{2q}(x)$.
Its Gram is $G_N=R_N^*M_NR_N$. For each $j\geq q-1$, let $p_{j+1-q}$
be monic and orthogonal to lower polynomials for the actual measure
$|\chi(S)|^2d\mu_x$. Its coefficients are the unique solution of the
positive lower-degree Gram system. Define
\[
 d_j=\chi p_{j+1-q},\qquad
 a_j=\langle d_j,d_j\rangle_x>0,\qquad
 \alpha_j=(\langle d_j,S^b\rangle_x)_{0\leq b<q}.
 \tag{BRU5}
\]
These $d_j$ are mutually orthogonal and span the successive original
relation spaces. Subtracting their orthogonal projections gives
\[
 \begin{aligned}
 R_N&=s-\sum_{j=q-1}^{N-1}d_j a_j^{-1}\alpha_j,\\
 G_N&=s^*M_{q-1}s-\sum_{j=q-1}^{N-1}a_j^{-1}\alpha_j^*\alpha_j,\\
 G_j&=G_{j+1}+a_j^{-1}\alpha_j^*\alpha_j,\\
 (\iota_i^N R_i)^*M_N(x)(\iota_j^N R_j)
    &=G_j\quad(q-1\leq i\leq j\leq N).
 \end{aligned}\tag{BRU6}
\]
Here $\iota_i^N:\mathcal P_i\hookrightarrow\mathcal P_N$ is the
literal coefficient inclusion. The common ambient space in the last
identity is any $\mathcal P_N$ with $N\geq j$. At the four endpoint
degrees we use $N=2q$, so this common metric is exactly $M_x$;
later products suppress only these specified inclusions.
Indeed each subtracted column is a relation, so $J_NR_N=I_E$.
Its orthogonality to the displayed basis of $\ker J_N$ proves
the minimum property by expanding the squared norm of every
other lift. Expansion gives the Gram and its rank-one difference.
Finally $R_i-R_j$ is a relation in $\mathcal P_j$, orthogonal
to $R_j$, proving the last identity. No limiting quotient is used.

At the four cutoffs $q-1,q,2q-1,2q$, every entry in BRU5--6 uses
only original moments through degree $4q$. Explicitly if
$d_j(S)=\sum_a d_{ja}S^a$, then
$\alpha_{jb}=\sum_a\overline{d_{ja}}
\int(\overline S)^aS^b\,d\mu_x$ and
$a_j=\sum_{a,b}\overline{d_{ja}}d_{jb}
\int(\overline S)^aS^b\,d\mu_x$.
For $\mu_{h,t}=\int y^t w_h(y)\,dy$, absolute convergence and
the ordinary multinomial expansion give the exact arithmetic input
\[
 \int y^d m_{h,k}(y)\,dy
 =\sum_{d_1+\cdots+d_k=d}\frac{d!}{d_1!\cdots d_k!}
                  \prod_{i=1}^k\mu_{h,d_i}.
 \tag{BRU7}
\]
Expanding $(k/2-iy)^a(k/2+iy)^b$ then computes the complex
source entries, with both signs and all powers of $k/2$ retained.

\subsection{The exact boundary update and its asymmetric overlap}

For the next-degree Gram $G=G_{j+1}$ put
\[
 \begin{gathered}
 H=I^*GI,\quad Q=(\Lambda G^{-1}\Lambda^*)^{-1},\quad
 S_0=G^{-1}\Lambda^*Q,\quad B_0=[I,S_0],\\
 b=\alpha_j I,\quad c=\alpha_jS_0,\quad a=a_j,\quad
 \beta=bH^{-1}b^*,\quad \nu=cQ^{-1}c^*.
 \end{gathered}\tag{BRU8}
\]
The section here is $S_0$, to avoid confusing it with the original
integration coordinate $S$. Direct multiplication gives
$\Lambda S_0=I_B$, $I^*GS_0=0$, and
$B_0^*GB_0=\operatorname{diag}(H,Q)$. Every vector has the unique
decomposition
\[
 z=I H^{-1}I^*Gz+S_0\Lambda z,
 \qquad G^{-1}=IH^{-1}I^*+S_0Q^{-1}S_0^*.
 \tag{BRU9}
\]
The first identity follows since $z-S_0\Lambda z\in K$, and its
kernel coefficient is recovered by $H^{-1}I^*G$. Multiplication
by $G$ proves the second. Consequently
$\alpha_jG^{-1}\alpha_j^*=\beta+\nu$, including all unit and
period dependence in $\Lambda$.

At degree $j$, the kernel Gram, boundary section, and boundary Gram
are exactly
\[
 \boxed{\begin{aligned}
 H_j&=H+b^*b/a,\\
 S_j&=S_0-IH^{-1}b^*c/(a+\beta),\\
 Q_j&=Q+c^*c/(a+\beta).
 \end{aligned}}\tag{BRU10}
\]
To prove the section formula, its image under $\Lambda$ is $I_B$.
Furthermore
\[
 I^*G_jS_j=b^*c/a-(H+b^*b/a)H^{-1}b^*c/(a+\beta)=0.
\]
Thus it is the unique orthogonal lift, by the direct decomposition
of $E$. Substituting it into $S_j^*G_jS_j$ gives the last line;
equivalently the lower block after eliminating the positive upper
block of
\[
 B_0^*G_jB_0=
 \begin{pmatrix}H+b^*b/a&b^*c/a\\c^*b/a&Q+c^*c/a\end{pmatrix}
\]
is $Q+c^*c/a-c^*b(H+b^*b/a)^{-1}b^*c/a^2$.
Multiplication verifies
$(H+b^*b/a)^{-1}b^*=aH^{-1}b^*/(a+\beta)$,
which gives exactly BRU10. All denominators are positive.

The determinant ratio is therefore the exact ordered product
\[
 \boxed{\frac{\det G_j}{\det G_{j+1}}
 =(1+\beta/a)\left(1+\frac{\nu}{a+\beta}\right)
 =1+\frac{\beta+\nu}{a}.}\tag{BRU11}
\]
Indeed the matrices $[I,S_j]$ and $[I,S_0]$ differ by a triangular
shear of determinant one. Their coordinate determinants cancel.
For any positive $Z$ and row $v$, the eigenvalues of
$Z^{-1/2}v^*vZ^{-1/2}$ are zero except for $vZ^{-1}v^*$;
this proves each rank-one determinant identity used here.

Let $\mathcal P_j^K$ and $\mathcal P_j^B$ be the original-source
orthogonal projections onto $R_jI(K)$ and $R_jS_j(B)$, respectively.
For successive degrees their four overlaps are
\[
 \boxed{\begin{aligned}
 \operatorname{Tr}(\mathcal P_j^K\mathcal P_{j+1}^K)
    &=d_K-\frac{\beta}{a+\beta},\\
 \operatorname{Tr}(\mathcal P_j^B\mathcal P_{j+1}^B)
    &=r-\frac{\nu}{a+\beta+\nu},\\
 \mathcal P_j^K\mathcal P_{j+1}^B&=0,\\
 \operatorname{Tr}(\mathcal P_j^B\mathcal P_{j+1}^K)
    &=\frac{\beta\nu}{(a+\beta)(a+\beta+\nu)}.
 \end{aligned}}\tag{BRU12}
\]
Here are proofs of every cross term. For $i\leq j$, BRU6 gives
the four cross Grams $H_j$, $Q_j$, zero, and $S_i^*G_jI$.
The zero follows from $G_jS_j=\Lambda^*Q_j$ and $\Lambda I=0$;
the boundary cross Gram is $S_i^*\Lambda^*Q_j=Q_j$.
For column maps $X,Y$ with positive Grams $H_X,H_Y$, substitution
of the projections $XH_X^{-1}X^*M$ and $YH_Y^{-1}Y^*M$ gives
trace $\operatorname{Tr}(H_X^{-1}X^*MYH_Y^{-1}Y^*MX)$.
This proves the incoming A10--12 on these actual source maps.
At consecutive degrees BRU10 gives
$X_{j,j+1}=S_j^*GI=-c^*b/(a+\beta)$.
The rank-one inverse identities give
\[
 bH_j^{-1}b^*=a\beta/(a+\beta),\qquad
 cQ_j^{-1}c^*=\nu(a+\beta)/(a+\beta+\nu).
\]
Substitution proves all four formulas in BRU12. Adding them proves
\[
 \operatorname{Tr}(G_j^{-1}G_{j+1})
       =q-\frac{\beta+\nu}{a+\beta+\nu}.
 \tag{BRU13}
\]
Thus neither squared contribution replaces the asymmetric cross
term: it is positive precisely when $\beta\nu>0$.

\subsection{All six cross-degree overlaps from complete relation blocks}

For $q-1\leq i<j\leq2q$, retain the ordered rows
$\mathcal A=(\alpha_i,\ldots,\alpha_{j-1})^{\mathsf T}$ and
$D=\operatorname{diag}(a_i,\ldots,a_{j-1})$. BRU6 gives
$G_i=G_j+\mathcal A^*D^{-1}\mathcal A$.
With $H=H_j,Q=Q_j,S_0=S_j$, put
\[
 b=\mathcal A I,\quad c=\mathcal A S_0,\quad
 Z=D+bH^{-1}b^*>0.
 \tag{BRU14}
\]
The complete block versions, retaining every correlation, are
\[
 \boxed{\begin{aligned}
 H_i&=H+b^*D^{-1}b,\\
 S_i&=S_0-IH^{-1}b^*Z^{-1}c,\\
 Q_i&=Q+c^*Z^{-1}c,\\
 X_{ij}&=-c^*Z^{-1}b,\\
 \frac{\det G_i}{\det G_j}
 &=\frac{\det H_i}{\det H}\frac{\det Q_i}{\det Q}.
 \end{aligned}}\tag{BRU15}
\]
For a direct verification,
$(H+b^*D^{-1}b)H^{-1}b^*=b^*D^{-1}Z$.
This identity gives $I^*G_iS_i=0$ and, after expansion of its
Gram, $Q_i=Q+c^*Z^{-1}c$. The displayed $X_{ij}$ follows by
multiplying $S_i^*$ by $G_jI$. The determinant argument is the
same shear of determinant one used above. This proves BRU15
without assuming that any of these blocks commute.

In particular the six pairs of four distinct endpoints are
computed by these specified finite matrices. Their overlaps are
\[
 \begin{aligned}
 o_{ij}^K&=\operatorname{Tr}(H_i^{-1}H_j),\qquad
 o_{ij}^B=\operatorname{Tr}(Q_i^{-1}Q_j),\\
 c_{ij}^{BK}&=\operatorname{Tr}
 (Q_i^{-1}c^*Z^{-1}bH_j^{-1}b^*Z^{-1}c),\qquad
 \mathcal P_i^K\mathcal P_j^B=0.
 \end{aligned}\tag{BRU16}
\]
The inverse $Z^{-1}$ retains all row correlations. Keeping only
its diagonal does not give this identity.

For $(N_0,N_1,N_2,N_3)=(q-1,q,2q-1,2q)$ and
$\epsilon=(1,1,-1,-1)$, define
$\mathsf S_Y=\sum_a\epsilon_a\mathcal P_{N_a}^Y$ for $Y=B,K$.
The complete squared norms calculated by BRU16 are
\[
 \begin{aligned}
 v_Y&=4\dim Y+2\sum_{a<b}\epsilon_a\epsilon_b o^Y_{N_aN_b},\\
 v_{\rm all}&=v_B+v_K+
             2\sum_{a<b}\epsilon_a\epsilon_b c^{BK}_{N_aN_b}
             =\operatorname{Tr}((\mathsf S_B+\mathsf S_K)^2).
 \end{aligned}\tag{BRU17}
\]
These equalities follow by expanding the squares, using idempotence
and the zero cross product in BRU16. The two remaining cross products
in the trace are equal by cyclicity, giving exactly the factor two.
No sign is assigned to their signed sum.

Write $T_x=M_x^{-1}(M_{\rm ar}-M_\Gamma)$ on the unchanged
$\mathcal P_{2q}$. Differentiating $J_NR_N=I$ makes $R_N'$ a
relation; its cross terms in $G_N'$ vanish by minimum orthogonality.
Thus $G_N'=R_N^*M_x'R_N$. Determinant differentiation gives
\[
 F'(x)=\operatorname{Tr}(T_x(\mathsf S_B+\mathsf S_K)),\qquad
 F=\log\frac{\det G_{q-1}\det G_q}{\det G_{2q-1}\det G_{2q}}.
 \tag{BRU18}
\]
Both $\mathsf S_Y$ have trace zero, since their four ranks are fixed.
Subtracting the midpoint of the extreme eigenvalues of $T_x$ in
this trace and applying Cauchy--Schwarz to the nonzero real
eigenvalues of the self-adjoint $\mathsf S_Y$ proves
\[
 |F'(x)|\leq\frac{\operatorname{spread}(T_x)}2
 \min\left\{\sqrt{\min(4r,2q+1)v_B}
            +\sqrt{\min(4d_K,2q+1)v_K},\,
                    \sqrt{(2q+1)v_{\rm all}}\right\}.
 \tag{BRU19}
\]
This rederives the incoming bound on its actual original domain
and supplies its complete finite relation evaluation BRU5--17.
The signed integral of BRU18 is exactly $F(1)-F(0)$; applying
BRU19 before integration is a bound for that same scalar.

\subsection{The exact two contributions to the arithmetic threshold}

For each $j$ evaluate BRU8 using $G_{j+1}(x)$, and denote the
result by $\beta_j(x),\nu_j(x)$. Define the two scalar sums
\[
 \begin{aligned}
 F_K(x)&=\sum_{j=q-1}^{2q-1}\omega_j
                \log(1+\beta_j(x)/a_j(x)),\\
 F_B(x)&=\sum_{j=q-1}^{2q-1}\omega_j
                \log\left(1+\frac{\nu_j(x)}{a_j(x)+\beta_j(x)}\right),\\
 \omega_{q-1}&=\omega_{2q-1}=1,\qquad
 \omega_j=2\quad(q\leq j\leq2q-2).
 \end{aligned}\tag{BRU20}
\]
Telescoping BRU11 proves, with every endpoint sign,
\[
 \boxed{F=F_K+F_B,\qquad
 \Delta_k^\Gamma=(F_K(1)-F_K(0))+(F_B(1)-F_B(0)).}
 \tag{BRU21}
\]
Each $F_Y(x)$ is nonnegative, since its individual arguments are
at least one. No sign follows for either endpoint difference.
These sums equal the signed logarithmic four-determinant ratios of
the canonical $H_N$ and $Q_N$, respectively. Those are the metrics
constructed on the same original source.
Changing a coefficient basis in $K$ or $B$ gives the same scalar
ratios by cancellation of its four congruence determinants.

On the domain $k\geq3$ of the inherited HC5 calculation, its original
constants and residual therefore have the exact receiving identity
\[
 F_K(1)+F_B(1)=4q\log(D_hk)+\mathcal R_k,
 \qquad\mathcal R_k\geq0.
 \tag{BRU22}
\]
Here $D_h$ and $\mathcal R_k$ retain the complete definitions HC3--4
in the dependency supplement; no phase, contraction or norm term
is reassigned to the boundary. BRU20 calculates what both terms
of this equality are. For $r=0$, $\nu_j=0$ and the entire scalar
is in $F_K$. For $r=q$, $\beta_j=0$ and it is in $F_B$.
For intermediate rank the cross response BRU12 is retained.
The new identities do not evaluate the displayed growing-family
moments or prove an upper bound opposing HC5.

\subsection{The same calculation on the original mixed graph source}

The formulas also act on the actual graph relation rows, with a
different stated Gram. Retain $s_i=1/2+it_i$, $\Sigma=\sum_i s_i$,
$\tau_k=\sum_i\Phi_h(s_i)$ and the original maps
\[
 \mathcal VP=\Bigl(\prod_i v_h(s_i)\Bigr)P(\Sigma),\quad
 F_b=\operatorname{rem}_{h(s_1),\ldots,h(s_k)}(\tau_k\Sigma^b),\quad
 L_{\rm obs}z=\Bigl(\prod_i v_h(s_i)\Bigr)\sum_bz_bF_b .
 \tag{BRU23}
\]
The graph lift of a polynomial $P$ is
$(\mathcal VP,\tau_k\mathcal VP-L_{\rm obs}J_NP)$ in the ordered
sum of two original Hilbert spaces, each with measure
$\prod_i dt_i/(2\pi)$. Define
\[
 r_k(y)=\int_{\sum t_i=y}(1+|\tau_k|^2)\prod_i w_h(t_i),\quad
 z_b(y)=\int_{\sum t_i=y}\overline{\tau_k}F_b\prod_iw_h(t_i).
 \tag{BRU24}
\]
The fibre integral eliminates the last variable with Lebesgue
measure, and for $k=1$ is evaluation. A relation $d=\chi Q$ has
graph $(\mathcal Vd,\tau_k\mathcal Vd)$, because $J_Nd=0$.
Its squared norm is exactly $\int|d(S)|^2r_k(y)\,dy$.
Its pairing with the graph of $S^b$ is exactly
$\int\overline{d(S)}(S^b r_k(y)-z_b(y))\,dy$, by expanding
the two ordered components. Thus the graph versions of BRU5 are
\[
 a_j^{\rm gr}=\int|d_j^{\rm gr}(S)|^2r_k(y)\,dy,\qquad
 (\alpha_j^{\rm gr})_b=
 \int\overline{d_j^{\rm gr}(S)}(S^b r_k(y)-z_b(y))\,dy,
 \tag{BRU25}
\]
where $d_j^{\rm gr}=\chi p_{j+1-q}^{\rm gr}$ and the monic
$p^{\rm gr}$ is orthogonal for $|\chi|^2r_k\,dy$.
These are precisely the complete PGRT/PGRC rows, with their cross
term. The same finite minimum proof gives
$\widehat G_j=\widehat G_{j+1}
+(a_j^{\rm gr})^{-1}(\alpha_j^{\rm gr})^*\alpha_j^{\rm gr}$.
Applying BRU8--17 and the two determinant sums of BRU20 to
these displayed rows and $\widehat G$
constructs $\widehat F_K,\widehat F_B$ and every graph boundary/kernel
overlap, on the same coefficient map $\Lambda$.
Here the determinant identity used from BRU21 is its first equality
$\widehat F=\widehat F_K+\widehat F_B$; the Gamma-to-arithmetic
derivative and endpoint-difference formulas keep their BRU4 path.

In particular the exact arithmetic return is
\[
 F(1)=\widehat F_K+\widehat F_B-\Gamma_k^{\rm gr},\qquad
 \Gamma_k^{\rm gr}=\sum_{a=0}^3\epsilon_a
                   \log\frac{\det\widehat G_{N_a}}{\det G_{N_a}(1)}.
 \tag{BRU26}
\]
This is determinant cancellation on the stated maps. The signed
augmentation is retained. The full PGRC inverse subtraction and
the mixed-source Gram $\Omega_k$ remain inside the graph metrics;
BRU26 does not give them the original moment norm.

Finally these are maps on every fixed coefficient face of the
original support diagram over $\tau$. An original relation
$\chi(\Sigma)Q(\Sigma)$ vanishes in $E_h^{\otimes k}$; successive
monic division constructs polynomials $Q_i$ with
$\chi(\Sigma)Q(\Sigma)=\sum_i h(s_i)Q_i(\mathbf s)$.
With the original source $h(D)F_h=\Theta\phi_*$ and
$D_i=-x_i\partial_{x_i}$, the primitive is
\[
 \sum_{i=1}^k(-1)^{i-1}Q_i(D_1,\ldots,D_k)
 (F_h^{\otimes(i-1)}\otimes\phi_*\otimes F_h^{\otimes(k-i)}).
 \tag{BRU27}
\]
The original source domains are stable under $D$: differentiation
preserves their Schwartz bounds and their vanishing moments, the
latter by integration by parts. Termwise differentiation of the
absolutely convergent theta sum gives $D\Theta=\Theta D$ there.
Consequently every displayed polynomial in the $D_i$ is an admitted
cochain operator and commutes with the tensor differential.
The tensor differential contributes the same sign and gives
$[\chi(\sum_i D_i)Q(\sum_iD_i)]F_h^{\otimes k}$ exactly.
The original Mellin map sends $F_h$ to $v_h$; integration by parts
on its original admitted domain sends $D_i$ to multiplication by
$s_i$, with its vanishing endpoint terms. Its ordered tensor map
therefore sends this boundary to $\mathcal V(\chi Q)$, with precisely
the measure $\prod_i dt_i/(2\pi)$ used above.
The two-leg inclusion has degree-zero
map $\phi\mapsto(\phi,0)$ and degree-one identity, so its differential
$\Theta(\phi-\mathcal F^{-1}\widehat\psi)$ preserves that equality.
All coefficient-face insertions commute because the displayed maps
are linear in each fixed coefficient and send zero to zero. On a
proper subcomplex $\mathsf C_\lambda$, when both compared top lifts
are admitted, the same boundary defines its exact residual in
$(d\mathsf C_{\rm full}^{k-1}\cap\mathsf C_\lambda^k)/
d\mathsf C_\lambda^{k-1}$. Its map to the kernel of the full top
cohomology inclusion and its inverse both use the same representative:
being zero in the full quotient means exactly being a full boundary.
This proves the kernel identification while retaining its primitive,
the two-leg diagonal, every independent mask, and the original outer
labels. None is quotiented out by the metric decomposition above.

For the graph relation the cochain realization is the ordered direct
sum of two copies of this tensor complex. Write $K_Q$ for BRU27 and
$T_\Phi=\sum_i\Phi_h(D_i)$. The proved commutation gives
$d(T_\Phi K_Q)=T_\Phi dK_Q$. Thus the actual primitive is
$(K_Q,T_\Phi K_Q)$; its differential, after the same original Mellin
map, is
$(\mathcal V(\chi Q),\tau_k\mathcal V(\chi Q))$, exactly the graph
relation used in BRU25. It has no $L_{\rm obs}$ term because
$J_N(\chi Q)=0$. Apply the direct sum of the two specified
$j_+^{\otimes k}$ maps to obtain its two-leg realization. This does
not identify the Euler operator on the Fourier leg with its operator
on the original leg. For a declared proper graph subcomplex the
residual is computed in this doubled full complex, on the inverse
image where both compared top lifts are admitted; the same
representative proof gives its kernel isomorphism.
